\section{Instalační návod}

\subsection{Požadavky}

Pro lokální spuštění projektu jsou potřeba následující nástroje:

\begin{itemize}
  \item Docker a~Docker Compose pro spuštění celé serverové a~webové vrstvy.
  \item Node.js 20 a~npm pro vývoj webové a~mobilní části.
  \item Bun pro přímé spuštění API serveru bez Dockeru.
  \item JDK~17 a~Android Studio (resp.\ Xcode na macOS) pro lokální sestavení mobilní aplikace.
\end{itemize}

\subsection{Spuštění celé webové a~serverové vrstvy přes Docker Compose}

Nejrychlejší lokální varianta je spuštění databáze, API serveru a~webového
frontendu společně pomocí Docker Compose.

\begin{enumerate}
  \item Vytvořte soubor \texttt{Server/.env} s~minimálně těmito proměnnými:
\end{enumerate}

\begin{verbatim}
PORT=3000
BETTER_AUTH_SECRET=change-me
POSTGRES_USER=lifemeter
POSTGRES_PASSWORD=lifemeter
POSTGRES_DB=lifemeter
\end{verbatim}

\begin{enumerate}
  \setcounter{enumi}{1}
  \item V~kořenovém adresáři repozitáře spusťte:
\end{enumerate}

\begin{verbatim}
docker compose up --build
\end{verbatim}

\begin{enumerate}
  \setcounter{enumi}{2}
  \item Po startu jsou k~dispozici tyto služby:
    \begin{itemize}
      \item API server: \texttt{http://localhost:3000/api}
      \item Webová prezentace a~download centrum: \texttt{http://localhost:3001}
      \item PostgreSQL: \texttt{localhost:5432}
    \end{itemize}
\end{enumerate}

\subsection{Spuštění API serveru bez Dockeru}

\begin{verbatim}
cd Server
bun install
PORT=3000 BETTER_AUTH_SECRET=change-me \
POSTGRES_USER=lifemeter POSTGRES_PASSWORD=lifemeter POSTGRES_DB=lifemeter \
bun run dev
\end{verbatim}

Swagger UI je následně dostupné na adrese \texttt{http://localhost:3000/api/ui}.

\subsection{Spuštění webové vrstvy bez Dockeru}

Při vývoji bez Dockeru je vhodné provozovat API server na portu~3000 a~Next.js
frontend na portu~3001.

\begin{enumerate}
  \item Vytvořte soubor \texttt{Web/.env.local} s~následujícím obsahem:
\end{enumerate}

\begin{verbatim}
API_URL=http://localhost:3000
NEXT_PUBLIC_API_URL=http://localhost:3000
APK_STORAGE_DIR=../releases/apk
IOS_STORAGE_DIR=../releases/ios
\end{verbatim}

\begin{enumerate}
  \setcounter{enumi}{1}
  \item Spusťte frontend:
\end{enumerate}

\begin{verbatim}
cd Web
npm ci
npm run dev -- --hostname 127.0.0.1 --port 3001
\end{verbatim}

\begin{enumerate}
  \setcounter{enumi}{2}
  \item Veřejné routy jsou pak dostupné na \texttt{http://localhost:3001/}
  a~\texttt{http://localhost:3001/downloads}, administrátorská část na
  \texttt{/admin}.
\end{enumerate}

\subsection{Spuštění mobilní aplikace}

\begin{verbatim}
cd App
npm ci
npm run android
\end{verbatim}

Na macOS lze obdobně použít \texttt{npm run ios}. Pro vývoj s~již
vygenerovaným dev klientem je možné využít také:

\begin{verbatim}
npm run start
\end{verbatim}

Publikační workflow \texttt{app-apk.yml} při změně adresáře \texttt{App/}
automaticky vytváří release APK, nahrává jej na produkční server do
adresáře \texttt{releases/apk} a~aktualizuje metadata využívaná webovou
stránkou \texttt{/downloads}.
